\chapter*{General Introduction}
\label{ch:intoGen}

Multi-agent reinforcement learning (MARL) is the application of reinforcement learning (RL) 
to multiple agent systems where the agents utilize their previous experiences to improve 
their performance through trial-and-error learning. Unlike traditional supervised learning 
approaches where learning occurs using an existing dataset, MARL allows agents to try out 
different possible actions and find alternatives to control and navigate in the environment 
without needing to have specific programming instructions to control actions in certain ways.

This research explores the application of MARL in developing a RL-based 3D physical soccer 
environment based on the MuJoCo physics engine. The 3D soccer simulation enables two teams 
of four players (2 vs 2) to play against each other without any set rules written into the 
code base of the soccer simulation application. The agents or players are required to learn 
how to control their bodies and coordinate their efforts in order to run to the ball, make 
passes to teammates, and kick the ball into the goal to score points. The task becomes 
progressively more complicated by using continuous action spaces, the interacting dynamic 
relationships between agents, as well as the rapidly changing environment due to the 
real-time actions of all five agents updating policy simultaneously.

In this project, we will focus our study on the theoretical underpinnings of modern 
reinforcement learning techniques by initially studying issues associated with discrete 
variables (tabular Q learning and discrete Q networks).

\section*{Organization of the report}
\label{sec:intro_org}

The remainder of this report will be organized as follows: Chapter~\ref{ch:into} describes 
the context of our host organization (ENSAM Casablanca), project scope and specific objectives; 
Chapter~\ref{ch:lit_rev} examines the mathematical basis for Reinforcement Learning, including 
Markov Decision Processes (MDPs), value based methods, and Continuous Actor-Critic algorithms; 
Chapter~\ref{ch:method} covers our methodology in detail, including environment setup, custom 
reward wrappers, and hyperparameter settings; Chapter~\ref{ch:results} provides our quantitative 
evaluation of the research data collected, with examples of both rewards and 2-dimensional 
spatial hexbin heat maps; Chapter~\ref{ch:evaluation} discusses cooperative behaviours among 
teams, team spreads, and physical limitations associated with cooperation; Chapter~\ref{ch:conclusion} 
wraps up this research project and provides thoughts on future directions for the project; 
and Chapter~\ref{ch:reflection} provides an individual engineering and learning reflection 
on the work done.
